\begin{proof}
We first check that \eqref{eq:p071-offcenter-quadratic} defines a generator.
On the first node axis,
\[
                         g_{ij}(r,0)=\frac12A_ir^2,
\]
which is independent of $j$.  On the second node axis,
\[
                         g_{ij}(0,s)=\frac12B_js^2,
\]
which is independent of $i$.  Thus the branch functions have the required
common smooth axis restrictions.  Every row and every column of $C$ sums
to zero.  Since
\[
 \partial_rg_{ij}(0,s)=C_{ij}s,
 \qquad
 \partial_sg_{ij}(r,0)=C_{ij}r,
\]
the two balance identities follow.

The Hessian of $g_{ij}$ is
\[
 H_{ij}=\begin{pmatrix}A_i&C_{ij}\\ C_{ij}&B_j\end{pmatrix}.
\]
All diagonal entries are positive, and direct calculation gives
\begin{equation}
 (A_iB_j-C_{ij}^2)_{ij}
 =\begin{pmatrix}
 82944&0&159744\\
 0&0&0\\
 0&0&49152
 \end{pmatrix}.
 \label{eq:p071-offcenter-determinants}
\end{equation}
Hence every $H_{ij}$ is positive semidefinite.  The nine restrictions are
jointly convex and $g\in\mathcal G_0$.

The weights in \eqref{eq:p071-offcenter-atoms} sum to one.  Their first
radial moments, grouped by the three branches of $S$, are
\begin{equation}
 (R_1,R_2,R_3)=
 \left(\frac7{172},\frac{25}{86},\frac7{172}\right).
 \label{eq:p071-offcenter-first-moments}
\end{equation}
Since $v_3=-v_1-v_2$, this gives
\[
 R_1v_1+R_2v_2+R_3v_3
  =(R_1-R_3)v_1+(R_2-R_3)v_2
  =\frac14v_2.
\]
The corresponding second radial moments are
\begin{equation}
 (S_1,S_2,S_3)=
 \left(\frac5{86},\frac{177}{86},\frac5{86}\right),
 \label{eq:p071-offcenter-second-moments}
\end{equation}
and therefore
\[
 S_1w_1+S_2w_2+S_3w_3=2w_2.
\]
This proves the ambient barycentre identity
\eqref{eq:p071-offcenter-barycentre}.

It remains to evaluate the quadratic.  In the order of
\eqref{eq:p071-offcenter-atoms}, its five values are
\begin{equation}
             50,\qquad \frac{121}{2},\qquad
             \frac{729}{32},\qquad \frac{841}{32},\qquad 2.
 \label{eq:p071-offcenter-values}
\end{equation}
At the proposed barycentre,
\[
 g(z_0)=g_{22}(1/4,2)=\frac{1681}{32}.
\]
Multiplying the values in \eqref{eq:p071-offcenter-values} by the five
weights and adding gives
\[
 \sum_a\lambda_ag(z_a)=\frac{57945}{1376}.
\]
Their difference is $7169/688>0$, proving
\eqref{eq:p071-offcenter-violation}.

The function $g$ already belongs to the stage-zero generator class and
hence to every generated sheaf under consideration.  Any comparison class
that imposed Jensen's inequality for all ambient barycentres would exclude
$g$.  Moreover, if the restriction of $g$ to these six points came from a
convex function on a convex ambient set containing them, ordinary Jensen
inequality would contradict \eqref{eq:p071-offcenter-violation}.  Thus the
unrestricted convex-roof route is unavailable, as claimed.
\end{proof}
